\section{Workspace.sty}\label{sec:ref-workspace}

Explicit handwriting whitespace, and the home of the overlay/solid
machinery. The concepts -- reserved space, the overlay contract, solid
mode -- are chapter~material (section~\ref{sec:overlay}); this is the
surface.

\subsection*{Commands}

\begin{center}
\begin{tabularx}{\linewidth}{@{}l Y@{}}
\cs{workspace} & 8 lines of handwriting room (the default).\\
\cs{workspace}\texttt{[12]} & that many \cs{HandwritingLine}s.\\
\cs{workspace}\texttt{[halfpage]} & half the text block, less one line --
  so a line at the page top sits vertically centred.\\
\cs{workspace}\texttt{[fullpage]} & the whole text block (text height, not
  paper).\\
\cs{workspace*} & fill to the bottom of the page (\cs{vfill}).\\
\cs{FreshPage} & collapsible page break; content flows onto the fresh
  page.\\
\cs{FreshPage}\texttt{[blank]} & a dedicated empty page; content starts
  the page after. The keyword is a visible reminder the blank is
  deliberate.\\
\cs{FreshPage}\texttt{[watermark]} & the blank page stamped with a faint
  diagonal \cs{BlankPageText}.\\
\cs{solidboxes} / \cs{overlayboxes} & the document-wide mode switch
  (section~\ref{sec:overlay}).\\
\end{tabularx}
\end{center}

A leaf command, never a wrapper: \cs{workspace} sits \emph{between}
content, so sources folded by content environments stay readable. Plain
\cs{newpage} is untouched -- use it for the rare break that must never
collapse.

\subsection*{Knobs}

\begin{center}
\begin{tabularx}{\linewidth}{@{}l Y@{}}
\cs{HandwritingLine} & the unit (0.9\,cm); every line-count workspace
  scales with it.\\
\cs{WSOverlayGap} & how far below the previous line an overlay box paints
  (\cs{medskipamount}), matched to the solid-mode top skip so the two
  modes look alike.\\
\cs{BlankPageText} & the watermark wording (``Blank page for extra
  work'').\\
\end{tabularx}
\end{center}

\subsection*{For package authors}

\cs{WSZeroHeight}\texttt{\{$\langle$a complete tcolorbox$\rangle$\}} is the
seam the box packages use: in overlay mode it emits the box at zero height
and depth, painting downward, with the following content placed exactly as
if nothing had been emitted; in solid mode it is a pass-through. Exercises
and InstructorNotes each carry an \cs{AtBeginDocument}
\cs{providecommand} identity fallback, so they work without this package
-- the narrow-coupling pattern discussed in
section~\ref{sec:architecture}. Two properties follow from the mechanism:
an overlay box never page-breaks (it is a fixed stamp), and one taller
than its workspace overprints -- the intended alarm.
